% !TEX root = ../main.tex
\chapter{Introduction}

\textbf{scala-man} is a single-player arcade maze game inspired by Pac-Man. Players navigate a
closed maze, collect collectibles, and avoid enemies to complete the selected game mode. It
is implemented in Scala 3, with a functional domain core and a ScalaFX user interface. The game
supports configurable maps, distinct enemy behaviours, and persistent results.

Each game starts with \textbf{three lives}. The player moves orthogonally between walkable cells,
while each enemy follows its own strategy. Clearing every \textbf{standard collectible} wins a
classic or timed game; survival has no collection objective and ends only when the player runs out of lives.
Meeting an enemy normally costs a life and returns the moving entities to their spawn positions.
Teleports create shortcuts by carrying an entity to the paired endpoint.

\section{Modes, score, and competition}

The same maze can be played in three modes, each with its own objective and ranking criterion:

\begin{itemize}
  \item \textbf{Classic} rewards standard collectibles, bonuses, defeated enemies, and lives left at the end; an
  invulnerability combo makes consecutive enemy defeats progressively more valuable.
  \item \textbf{Timed} applies the classic objective under a time limit. The seconds remaining are shown as a
  potential score, but are assigned and used to rank a run only when the maze is cleared; losing every life or
  reaching the limit yields zero points.
  \item \textbf{Survival} has no victory condition. Every \textbf{thirty seconds} raises enemy speed by
  \textbf{20\%}, up to the player's base speed, and the score records the difficulty waves survived.
\end{itemize}

Results are stored in a \textbf{leaderboard for each map-mode combination}, ordered by descending score.
Only comparable runs are therefore ranked together.

\section{Bonuses}

In addition to standard collectibles, mazes can contain any number of bonus items of two \textbf{types}:
\textbf{Invulnerability} and \textbf{Slowdown}.

\begin{itemize}
  \item \textbf{Invulnerability} lasts \textbf{eight seconds}. During this interval, meeting an enemy does not
  cost a life; instead, the enemy is defeated. Consecutive enemy defeats during the same effect also build
  the score combo in Classic mode.
  \item \textbf{Slowdown} lasts \textbf{five seconds}. It halves the speed of every enemy while leaving the
  player at normal speed, creating time to escape, collect a standard collectible, or choose a safer route.
\end{itemize}

The player can use either effect defensively to survive a dangerous crossing or offensively to improve a score.

\section{Enemy tactics}

The maze contains three kinds of enemy:

\begin{itemize}
  \item the \textbf{Hunter}, which follows a shortest path towards the player's current position;
  \item the \textbf{Anticipator}, which predicts where the player is heading and attempts to intercept them;
  \item the \textbf{Patroller}, which ignores the player and repeatedly follows a cyclic route through the
  walkable cells nearest the maze corners.
\end{itemize}

\section{Maps, persistence, and project goals}

Players can choose a \textbf{supplied maze} or provide their own through the
\href{https://github.com/PPS-25/PPS-25-scala-man/blob/dev/docs/map-format.md}{\textbf{ASCII map-format document} in the project repository}.
Before play, a map is
\textbf{parsed and validated}: it must be non-empty and rectangular, use only supported symbols, have a closed border,
exactly one player spawn, at least one standard collectible and one enemy, correctly paired teleports, and every
standard collectible and enemy reachable from the spawn. \textbf{Invalid maps} are rejected with a clear explanation rather
than causing the game to crash.

A paused game can be \textbf{saved} from the in-game menu and \textbf{resumed} later with its maze, positions, remaining collectibles,
effects, score, clock, and mode restored. A result is recorded only when a game reaches a \textbf{terminal outcome}.

Patrolling enemies, additional game modes, score leaderboards, and game persistence were planned as
\textbf{optional extensions} in the project proposal; all were completed and integrated into the final application.

Beyond the playable application, scala-man pursues three design goals developed in the following chapters:
\textbf{extensible enemy behaviour}, \textbf{extensible game modes}, and a \textbf{testable domain model}
that remains independent from the graphical interface. The report then presents the development process,
requirements, architecture, detailed design, implementation, testing, and retrospective.
